% Example "pleiades": the problem, the seven orbits with their close
% encounters, and the convergence and invariant tests. Every number comes from
% data/ (parameters.tex, results.tex, orbits.dat, encounter.dat,
% encounters.dat, history_<method>.dat, energy_<method>.dat,
% error1_<method>.dat, error_<method>.dat), written by
% examples/pleiades/main.f90: run `fpm run --example pleiades` from the project
% root, then build this file with lualatex in examples/pleiades/.
\documentclass[11pt,a4paper]{article}
\usepackage[margin=2.2cm]{geometry}
\usepackage{amsmath,booktabs,array}
\usepackage{pgfplots}
\usepgfplotslibrary{groupplots,fillbetween}
\pgfplotsset{compat=1.18}
\pagestyle{empty}

% Okabe-Ito colour-blind-safe palette, one colour per body.
\definecolor{b1}{HTML}{E69F00}  % orange
\definecolor{b2}{HTML}{56B4E9}  % sky blue
\definecolor{b3}{HTML}{009E73}  % bluish green
\definecolor{b4}{HTML}{0072B2}  % blue
\definecolor{b5}{HTML}{D55E00}  % vermilion
\definecolor{b6}{HTML}{CC79A7}  % reddish purple
\definecolor{b7}{HTML}{444444}  % dark grey
% Method colours, the same in every example.
\colorlet{m-euler}{black}
\colorlet{m-order2}{orange!85!black}
\colorlet{m-order3}{green!50!black}
\colorlet{m-rk4}{blue!80!black}
\colorlet{m-dp54}{red!75!black}

\begin{document}
\section*{Example: Pleiades, seven bodies}

\subsection*{1. Problem}
From the SciML problem library (\texttt{prob\_ode\_pleiades}), after Hairer,
N{\o}rsett and Wanner (\emph{Solving ODEs I}, section II.10): seven stars in
a plane, star $j$ of mass $m_j = j$, attracting each other by gravity
($G = 1$),
\begin{equation*}
  \ddot x_i = \sum_{j \ne i} m_j\,\frac{x_j - x_i}{r_{ij}^3}, \qquad
  \ddot y_i = \sum_{j \ne i} m_j\,\frac{y_j - y_i}{r_{ij}^3}, \qquad
  r_{ij} = \sqrt{(x_i - x_j)^2 + (y_i - y_j)^2},
\end{equation*}
a system of 28 first-order equations for $0 \le t \le 3$. There is no closed
form; the stars pass close to each other several times, where the forces
become large and change fast. The reference is a \texttt{dp54} run with
over a million steps. Three quantities are conserved: the total momentum
$P = \sum m_i \dot q_i$, the angular momentum
$L = \sum m_i (x_i\dot y_i - y_i\dot x_i)$ and the energy
$E = \tfrac12\sum m_i|\dot q_i|^2 - \sum_{i<j} m_i m_j / r_{ij}$.
\begin{center}
\small
% Written by examples/pleiades/main.f90: do not edit.
\begin{tabular}{l>{\raggedright\arraybackslash}p{12.4cm}}
\toprule
masses & $m_j = j$, $j = 1, \dots, 7$; $G = 1$ \\
positions & $(3,3), (3,-3), (-1,2), (-3,0), (2,0), (-2,-4), (2,4)$ \\
velocities & zero except $\dot x_6 = 1.75$, $\dot x_7 = -1.5$, $\dot y_4 = -1.25$, $\dot y_5 = 1$ \\
interval & $0 \le t \le 3$ \\
reference & \texttt{dp54}, 1572864 steps; $2.11\times10^{-10}$ from the run with half the steps \\
convergence & error $\max_i |u_i - u_i^{\mathrm{ref}}|$ at $t = 1$, $3$; 384, 768, \dots, 196608 steps \\
observed order & two finest levels above $3.00\times10^{-11}$ ($t = 1$), $1.00\times10^{-8}$ ($t = 3$); kept if within 0.3 of the next slope \\
invariants, history & 3072 steps ($\Delta t = 0.001$); $E_0 = -45.9525$, $L_0 = 109$, $P_0 = (0, 0)$ \\
methods & \texttt{rk3ws}, \texttt{rk4}, \texttt{dp54} \\
interface & plain arrays, \texttt{rk\_array\_scheme\_t}, 28 unknowns \\
\bottomrule
\end{tabular}

\end{center}

\newpage
\subsection*{2. Solutions}
\begin{center}
\begin{tikzpicture}
  % ---- (a) the seven orbits ----
  \begin{axis}[
      name=orbits,                                         % anchor for the panels on the right
      width=9.4cm, height=13.4cm,                          % tall: the orbits span y = -4 ... 5.2
      axis equal image,                                    % true shapes
      xmin=-3.9, xmax=3.9, ymin=-4.6, ymax=5.8,            % all orbits, with margins
      xlabel={$x$}, ylabel={$y$},                          % the plane of motion
      title={(a) reference orbits, $0 \le t \le 3$},       % panel title
      grid=major, grid style={black!8},                    % faint grid
      axis background/.style={fill=black!2},               % very light background
    ]
    \foreach \k/\w in {1/0.6, 2/0.75, 3/0.9, 4/1.05, 5/1.2, 6/1.35, 7/1.5} {
      \edef\temp{\noexpand\addplot[color=b\k, line width=\w pt, line join=round]
        table[x=x\k, y=y\k] {data/orbits.dat};}          % orbit of body k, width grows with its mass
      \temp
      \edef\temp{\noexpand\addplot[only marks, mark=*, mark size=2.2pt, color=b\k]
        table[x=x\k, y=y\k, restrict expr to domain={\noexpand\thisrow{t}}{2.999:3.001}] {data/orbits.dat};}
      \temp                                                % filled dot: position at t = 3
    }
    % starting positions, hollow, numbered by body (mass)
    \foreach \k/\x/\y/\pos in {1/3/3/east, 2/3/-3/east, 3/-1/2/north east, 4/-3/0/west,
                               5/2/0/south east, 6/-2/-4/south, 7/2/4/north} {
      \edef\temp{\noexpand\node[circle, draw=b\k, fill=white, line width=0.9pt, inner sep=1.6pt]
        at (axis cs:\x,\y) {};}                            % hollow circle: position at t = 0
      \temp
      \edef\temp{\noexpand\node[font=\noexpand\scriptsize\noexpand\bfseries, text=b\k, anchor=\pos, inner sep=3pt]
        at (axis cs:\x,\y) {\k};}                          % body number = mass
      \temp
    }
  \end{axis}
  % ---- (b) closest distance in time ----
  \begin{axis}[
      name=distance,                                       % anchor for panel (c)
      at={($(orbits.north east)+(1.6cm,-0.55cm)$)}, anchor=north west, % to the right of (a)
      width=7.2cm, height=5.6cm,                           % short panel
      xmin=0, xmax=3, ymode=log, ymin=0.02, ymax=6,        % distances from 0.03 to 5
      ylabel={$\min_{i<j} r_{ij}$},                        % smallest distance between two bodies
      xticklabels={},                                      % shares the t axis with (c)
      title={(b) closest pair},                            % panel title
      grid=major, grid style={black!8},                    % faint grid
      axis background/.style={fill=black!2},               % very light background
    ]
    \addplot[name path=d, black, thick] table[x=t, y=distance] {data/encounter.dat};   % smallest distance
    \path[name path=floor] (axis cs:0,0.02) -- (axis cs:3,0.02);                      % bottom of the axis
    \addplot[black!12] fill between[of=d and floor];                                  % shade under the curve
    \addplot[only marks, mark=*, mark size=1.8pt, red!75!black,
             nodes near coords, point meta=explicit symbolic,                          % label each encounter
             every node near coord/.style={font=\tiny, rotate=90, anchor=west, xshift=2pt, text=red!60!black, fill=white, fill opacity=0.8, text opacity=1, inner sep=1pt}]
      table[x=t, y=distance, meta=pair] {data/encounters.dat};                         % close encounters, by pair
  \end{axis}
  % ---- (c) position error of the methods ----
  \begin{axis}[
      at={($(distance.south west)+(0,-1.1cm)$)}, anchor=north west, % below (b), room for its title
      width=7.2cm, height=7.4cm,                           % taller panel
      xmin=0, xmax=3, ymode=log, ymin=1e-12, ymax=10,      % errors from dp54 to order one
      xlabel={$t$}, ylabel={position error},               % error against the reference
      title={(c) error, 3072 steps},                       % panel title
      grid=major, grid style={black!8},                    % faint grid
      axis background/.style={fill=black!2},               % very light background
      legend style={font=\tiny, at={(0.98,0.03)}, anchor=south east, fill=white, fill opacity=0.85,
                    text opacity=1, draw=black!30},        % compact legend, in the empty bottom-right corner
      legend cell align=left,
    ]
    \addplot[ycomb, red!25, thin, forget plot]
      table[x=t, y expr=1e3] {data/encounters.dat};       % vertical lines at the encounters
    \addplot[m-order3, thick, dashed] table[x=t, y=error] {data/history_rk3ws.dat};
    \addlegendentry{\texttt{rk3ws}}
    \addplot[m-rk4, thick]    table[x=t, y=error] {data/history_rk4.dat};
    \addlegendentry{\texttt{rk4}}
    \addplot[m-dp54, thick]   table[x=t, y=error] {data/history_dp54.dat};
    \addlegendentry{\texttt{dp54}}
  \end{axis}
\end{tikzpicture}
\end{center}
(a) The seven orbits of the reference: hollow circles are the starting
positions, numbered by body (the number is also the mass, and the line
width grows with it), filled dots the positions at $t = 3$. (b) The distance
between the closest two stars: seven close encounters (red, labelled by the
pair), the closest between stars 1 and 7 at $t \approx 1.68$, at a distance
of $0.034$. (c) The error of each method at 3072 steps against the
reference: it jumps at every encounter (light red lines), by orders of
magnitude for the low-order methods; before the first encounter
($t \approx 1.23$) the errors are small and grow smoothly.

\newpage
\subsection*{3. Tests: convergence and the invariants}

\paragraph{Prediction.} (i) Each method converges at its general order;
close encounters add large higher-order terms, so the asymptotic range
starts at smaller steps after them. \texttt{rk3ws} is a special case: its
only third-order defect is the local error $-\tfrac{h^3}{24}f''(f,f)$ (here
the velocity error $-\tfrac{h^3}{24}g''(q)[\dot q, \dot q]$), which equals
$h^3\bigl(\tfrac{d}{dt}X - f'X\bigr)$ for $X = -\tfrac1{24}\ddot u$ on any
autonomous problem; its second-order global error,
$h^2\bigl(X(u(t)) - \Phi(t,0)X(u_0)\bigr)$, therefore does not accumulate
like the third-order one (effective order 3; the nonlinear chemistry example
removes it by pre- and post-processing). (ii) The total momentum $P$ is a linear
invariant: every Runge--Kutta method keeps it to round-off. $L$ (quadratic)
and $E$ are kept by none of these methods.

\paragraph{Result.}
\begin{center}
\small
% Written by examples/pleiades/main.f90: do not edit.
\begin{tabular}{lccccccc}
\toprule
& & \multicolumn{2}{c}{observed order} & \multicolumn{3}{c}{largest drift, 3072 steps} \\
\cmidrule(lr){3-4}\cmidrule(lr){5-7}
name & order & $t = 1$ & $t = 3$ & momentum $|P - P_0|$ & angular $|L - L_0|$ & energy $|E - E_0|$ \\
\midrule
\texttt{rk3ws} & 2 & 1.99 & 3.00 & $5.06\times10^{-14}$ & $1.40\times10^{-2}$ & $5.30\times10^{0}$ \\
\texttt{rk4} & 4 & 4.02 & 4.21 & $1.38\times10^{-13}$ & $2.73\times10^{-4}$ & $2.10\times10^{-1}$ \\
\texttt{dp54} & 5 & --- & --- & $1.36\times10^{-13}$ & $4.11\times10^{-5}$ & $2.60\times10^{-2}$ \\
\bottomrule
\end{tabular}

\end{center}

\begin{center}
\begin{tikzpicture}
  \begin{groupplot}[
      group style={group size=3 by 1, horizontal sep=1.3cm, y descriptions at=edge left}, % three panels
      width=5.6cm, height=6cm,                             % size of each panel
      grid=major, grid style={black!8},                    % faint grid
      axis background/.style={fill=black!2},               % very light background
      every axis plot/.append style={thick, mark size=1.2pt}, % same weight for all curves
    ]
    % ---- (i) error at t = 1 ----
    \nextgroupplot[
      xmode=log, ymode=log, ymin=1e-13, ymax=1e3,          % shared error range
      title={(i) error at $t = 1$}, xlabel={$\Delta t$}, ylabel={error},
      legend to name=pltestlegend, legend columns=7,       % legend below the panels
      legend style={font=\small, /tikz/every even column/.append style={column sep=0.25cm}},
    ]
    \addplot[m-order3, dashed, mark=triangle] table[x=dt, y=error] {data/error1_rk3ws.dat};
    \addlegendentry{\texttt{rk3ws}}
    \addplot[m-rk4, mark=diamond*]           table[x=dt, y=error] {data/error1_rk4.dat};
    \addlegendentry{\texttt{rk4}}
    \addplot[m-dp54, mark=pentagon*]         table[x=dt, y=error] {data/error1_dp54.dat};
    \addlegendentry{\texttt{dp54}}
    % ---- (ii) error at t = 3 ----
    \nextgroupplot[
      xmode=log, ymode=log, ymin=1e-13, ymax=1e3,          % shared error range
      title={(ii) error at $t = 3$}, xlabel={$\Delta t$},
    ]
    \addplot[m-order3, dashed, mark=triangle] table[x=dt, y=error] {data/error_rk3ws.dat};
    \addplot[m-rk4, mark=diamond*]           table[x=dt, y=error] {data/error_rk4.dat};
    \addplot[m-dp54, mark=pentagon*]         table[x=dt, y=error] {data/error_dp54.dat};
    % ---- (iii) energy drift in time ----
    \nextgroupplot[
      ymode=log, ymin=1e-13, ymax=1e3,                     % shared range
      title={(iii) $|E - E_0|$, 3072 steps}, xlabel={$t$}, xmin=0, xmax=3,
    ]
    \addplot[ycomb, red!25, thin] table[x=t, y expr=1e3] {data/encounters.dat};  % encounters
    \addplot[m-order3, dashed, no marks] table[x=t, y=drift] {data/energy_rk3ws.dat};
    \addplot[m-rk4, no marks]      table[x=t, y=drift] {data/energy_rk4.dat};
    \addplot[m-dp54, no marks]     table[x=t, y=drift] {data/energy_dp54.dat};
  \end{groupplot}
  \node[below=0.2cm] at (current bounding box.south)        % legend centred under the panels
    {\pgfplotslegendfromname{pltestlegend}};
\end{tikzpicture}
\end{center}
(i) Before the first encounter ($t = 1$) every method shows its general
order, \texttt{rk3ws} 2; \texttt{dp54} is at round-off
already at the largest step (---). (ii) After the encounters ($t = 3$) the
errors are far larger (at $\Delta t \approx 10^{-3}$, up to
$6 \times 10^{10}$ times for \texttt{rk4}) and the asymptotic range starts at
much smaller steps: \texttt{dp54} never settles in this range
(---); \texttt{rk3ws} now shows order 3 over five levels, because its
third-order error, amplified at each encounter, outweighs its small
second-order term at every step tested. (iii) The energy error jumps at the
encounters and stays; the total momentum, a linear invariant, is kept to
$10^{-13}$ by every method (table).
\end{document}
